% 中文
\cabstract{
本实验基于 Oberon-0 语言定义 \cite{lab05}，设计并实现了一个面向 Oberon-0 语言的逆向工程工具 ROSE（Reverse Oberon Software Engineering）。实验分为四个子实验：实验一通过编写正确的 Oberon-0 源程序及其变异程序，深入理解 Oberon-0 语言的词法、语法和语义定义，讨论了保留字与关键字的区别以及文法的二义性问题；实验二利用词法分析程序自动生成工具 JFlex，编写 Oberon-0 语言的词法规则正则定义式，自动生成词法分析程序并对三种 lex 族工具进行比较；实验三使用语法分析程序自动生成工具 JavaCUP，编写 Oberon-0 语言的语法规则和语义动作，自动生成支持调用图绘制的语法分析与语法制导翻译程序；实验四手工设计 Oberon-0 语言的翻译模式，采用递归下降预测分析技术编写语法分析程序，实现流程图自动生成，并对自顶向下与自底向上两种分析策略进行了深入的比较分析。本实验通过 ROSE 工具的完整设计与实现，帮助学生深入理解和牢固掌握编译技术中的词法分析、语法分析、语法制导翻译以及自动生成工具等重要环节。
% TODO：等待补充结果摘要
}

\ckeywords{编译原理；Oberon-0；逆向工程；词法分析；语法分析；JFlex；JavaCUP；递归下降；ROSE}